% -------------------------------------------------
% VERSION TRACKER — No modificar este bloque
% Versión:    Tesis Humanizada
% Archivo:    Capitulo1Introduccion/Metodologia.tex
% Última mod:  2026-05-08T08:18:14
% Git commit:  cc9c766f @ main
% Audiencia:   general
% Verificar antes de editar: bash check_tesis_version.sh overleaf_tesis_humanizado/Capitulo1Introduccion/Metodologia.tex
% -------------------------------------------------


\section{Metodología}

\subsection{Tipo de Investigación}

Este proyecto es de \textbf{investigación aplicada}: toma el conocimiento de la ingeniería de sistemas y lo usa para resolver un problema real del entorno empresarial colombiano. También es \textbf{investigación proyectiva}, porque propone y construye una solución de software funcional —no se queda en la teoría, sino que entrega un producto que funciona. El proyecto sigue los lineamientos de la Resolución núm. 137 de la Facultad de Ingeniería de la Universidad del Valle para trabajos de grado en la modalidad de Trabajo de Investigación \cite{univalle_resolucion137}.

\subsection{Enfoque Metodológico: SCRUM}

Para gestionar el desarrollo del sistema se planteó \textbf{SCRUM}, una metodología ágil de trabajo que organiza el proyecto en ciclos cortos llamados \textit{sprints}. ¿Por qué SCRUM y no una metodología tradicional? La respuesta es sencilla: cuando se construye software para un usuario real (en este caso, una panificadora), es imposible anticipar todas las necesidades desde el primer día. SCRUM permite ir construyendo, mostrando avances, recibiendo retroalimentación y ajustando el rumbo cada dos semanas —en lugar de esperar meses para descubrir que el sistema no sirve.

SCRUM define tres roles principales. El plan los adaptó de la siguiente manera:

\begin{itemize}
    \item \textbf{Product Owner}: Rol que representa al cliente y define qué debe hacer el sistema. En Pymetory, esta función la cumple el director del trabajo de grado (Prof. Héctor Fabio Ocampo), quien definió las reglas de negocio del proyecto y revisa los avances en las reuniones de seguimiento.
    \item \textbf{Equipo de Desarrollo}: El autor del proyecto (Germán David Murillas Mondragón), que construye el sistema de manera unipersonal.
    \item \textbf{Scrum Master}: Rol que facilita el proceso y elimina obstáculos. También lo asume el autor.
\end{itemize}

\subsection{Cómo se organizó el trabajo}

El plan dividió el proyecto en 6 sprints de dos semanas cada uno —12 semanas de desarrollo activo— precedidos por 4 semanas de levantamiento de requisitos y diseño. Cada sprint contemplaba cuatro momentos:

\begin{itemize}
    \item \textbf{Product Backlog (lista de pendientes)}: Listado priorizado de las funcionalidades del sistema, gestionado en GitHub Projects.
    \item \textbf{Sprint Planning (planeación del ciclo)}: Al inicio de cada sprint se eligen las tareas del backlog para las dos semanas siguientes.
    \item \textbf{Sprint Review (revisión del avance)}: Al final de cada sprint se muestra al director un avance de software funcionando, para recibir retroalimentación y ajustar prioridades.
    \item \textbf{Sprint Retrospective (lecciones aprendidas)}: Después de la revisión se identifica qué funcionó, qué mejorar y qué acción tomar en el siguiente sprint.
\end{itemize}

El plan de sprints se detalla en el Capítulo 4. Con un solo desarrollador, en la práctica el trabajo se gestionó como un flujo continuo sobre un tablero Kanban en GitHub Projects, con criterios de aceptación y una \textit{Definition of Done}; esa ejecución real se describe en el Capítulo \ref{chap:gestion}.

\newpage
